\documentclass[11pt,a4paper,landscape]{article}
\usepackage[margin=16mm]{geometry}
\usepackage{qpp-tikz}
\usepackage{quantikz}
\usepackage[hidelinks]{hyperref}
\setlength{\parindent}{0pt}
\setlength{\parskip}{7pt}
\begin{document}
\section*{One operator, two syntaxes}
These trees represent the same two-qubit operator, of dimension $4$.
The UOp is produced directly by \texttt{fromQOp}, with no additional normalization.
Operands are ordered from left to right; at a composition node, the rightmost acts first.
\begin{center}
\begin{minipage}[t]{.24\textwidth}\centering
\textbf{QOp (source)}\par\medskip\input{qop-example.tex}
\end{minipage}\hfill
\begin{minipage}[t]{.73\textwidth}\centering
\textbf{UOp (translation)}\par\medskip\input{uop-example.tex}
\end{minipage}
\end{center}
Both roots are annotated $:\,2$. The UOp root's subscript in $\circ_4$ records
the dimension, while its gray annotation records the qubit count.
Lowering expands $H$ into rotations and a scalar phase, and replaces control by
a direct sum with an identity block.

In both syntaxes, $R(t)$ has its axis $A$ as a child and means $e^{-i\pi tA/2}$.
Identity subscripts denote dimensions. Permutations retain the input dimensions
and index list: $P_\otimes(ds;p)$ or $P_\oplus(ds;p)$, where output position $j$
selects input $p_j$.

\newpage
\section*{Empty spaces, permutations, and adjoints}
Empty tensor and direct sum have different dimensions; their labels keep this distinction visible.
\begin{center}
\input{uop-empty-tensor.tex}\qquad\input{uop-empty-sum.tex}\qquad\input{uop-identity.tex}
\par\bigskip\input{qop-permute.tex}\qquad\input{uop-permute.tex}
\par\bigskip
\begin{minipage}[t]{.28\textwidth}\centering
\textbf{Adjoint}\par\medskip
\input{qop-adjoint.tex}
\end{minipage}
\begin{minipage}[t]{.38\textwidth}\centering
\textbf{Binary composition}\par\medskip
\input{qop-gates.tex}
\end{minipage}
\begin{minipage}[t]{.28\textwidth}\centering
\textbf{Scalar composition}\par\medskip
\input{qop-scalars.tex}
\end{minipage}
\end{center}
The printer preserves exact rational angles, arbitrary dimensions, nested
associations, and singleton lists. It neither normalizes nor type-checks its input.

\newpage
\section*{Qubit counts}
The annotation $:\,n$ outside each outline gives the number of qubits.
Scalars act on zero qubits; tensor adds counts, control adds one, and a direct
sum of equal-sized blocks adds one. Composition and rotation retain the width.
\begin{center}
\begin{minipage}[t]{.48\textwidth}\centering
\textbf{QOp}\par\medskip\input{qop-annotated.tex}
\end{minipage}\hfill
\begin{minipage}[t]{.48\textwidth}\centering
\textbf{Direct sum}\par\medskip\input{qop-sum-annotated.tex}
\end{minipage}
\end{center}
\begin{center}
\begin{minipage}[t]{.48\textwidth}\centering
\textbf{Independent UOp example (dimension 3)}\par\medskip\input{uop-annotated.tex}
\end{minipage}\hfill
\begin{minipage}[t]{.48\textwidth}\centering
\textbf{Scalar and zero-dimensional spaces}\par\medskip\input{uop-empty-annotated.tex}
\end{minipage}
\end{center}
For UOp, dimension $2^n$ gives count $n$. Other dimensions are marked $d=\cdots$;
in particular, $d=0$ is not a zero-qubit space. Counts assume well-formed input.
The independent UOp example combines blocks of dimensions $1$ and $2$, giving $3$.
Counts appear by default. Use \texttt{showTikZQOpWithoutQubits} or
\texttt{showTikZUOpWithoutQubits} to omit the outside annotations.

\newpage
\section*{A small QFT: QOp}
This is \texttt{qft 2}, retaining its recursive construction, identity padding,
and binary associations. Following the edges makes each operand's scope explicit.
\begin{center}\input{qft2-qop-annotated.tex}\end{center}

\newpage
\section*{A small QFT: normalized UOp}
The same \texttt{qft 2} is lowered and normalized explicitly before printing.
Its outer composition has eight children, read from right to left.
\begin{center}\input{qft2-uop-without-qubits.tex}\end{center}
Each Hadamard becomes three rotation factors with its phase retained.
Printing only traverses syntax; no operator matrix is constructed.
This wide example uses \texttt{showTikZUOpWithoutQubits} to keep the tree compact.

\newpage
\section*{Flat circuits}
\texttt{QPP.QCirc} reads an operator as gates on wires, and \texttt{showQuantikz}
draws the result with the \href{https://ctan.org/pkg/quantikz}{quantikz} package.
Gates are placed in columns as soon as possible without two gates covering a common
row. Controls link to their targets; a box on non-consecutive wires spans them, numbers
its inputs by local qubit, and lets the wires it does not act on pass through.
A gate on no wires is a scalar, emitted as a comment.
\begin{center}
\input{circuit-bell.tex}\qquad\input{circuit-elements.tex}
\end{center}
\texttt{qft 3} as extracted by \texttt{fromQOp}, then \texttt{qft 4} after \texttt{simplify}:
\begin{center}
\input{qft3-circuit.tex}
\par\bigskip
\input{qft4-circuit-simplified.tex}
\end{center}
\texttt{showQCircDiagram} draws the same layout in line-drawing characters for the terminal.

\newpage
\section*{Restyle the same tree}
The diagram below reuses \texttt{qop-example.tex} unchanged.
TikZ controls appearance; Forest controls spacing and tree layout.
\begin{verbatim}
\tikzset{
  qpp/node/.append style={rounded corners=0pt, line width=.8pt},
  qpp/gate/.append style={fill=yellow!20},
  qpp/edge/.append style={thick, draw=blue!60!black}
}
\forestset{qpp tree/.append style={for tree={l sep=5mm, s sep=7mm, fork sep=2.5mm}}}
\end{verbatim}
\begingroup
\tikzset{
  qpp/node/.append style={rounded corners=0pt, line width=.8pt},
  qpp/gate/.append style={fill=yellow!20},
  qpp/edge/.append style={thick, draw=blue!60!black}
}
\forestset{qpp tree/.append style={for tree={l sep=5mm, s sep=7mm, fork sep=2.5mm}}}
\begin{center}\input{qop-example.tex}\end{center}
\endgroup
\small The Haskell printer translates each operator into \texttt{Data.Tree},
then emits \href{https://ctan.org/pkg/forest}{Forest}'s bracket notation.
All layout and styling live in \texttt{qpp-tikz.sty}; no coordinates or drawing macros
are generated.
\end{document}
